% !TEX root = ../thesis.tex
% ==============================================================================
%  CHAPTER 5: EXPERIMENTS & RESULTS (فصل پنجم: آزمایش‌ها و نتایج تجربی) (chapters/chapter5_experiments.tex)
% ==============================================================================
%  This chapter presents empirical validation, benchmark evaluations,
%  interventional ablation studies, efficiency metrics, and more.
% ==============================================================================

\chapter{آزمایش‌ها و ارزیابی نتایج}\label{chap:chapter5_experiments}

فصل پنجم پایان‌نامه به ارزیابی تجربی و تحلیل نتایج اختصاص دارد. در این فصل، روش پیشنهادی در شرایط عملیاتی یا شبیه‌سازی‌های استاندارد به بوته‌ی آزمایش گذاشته شده و عملکرد آن در مقایسه با روش‌های شاخص موجود ارزیابی می‌گردد. هدف از این آزمایش‌ها، اثبات کارآمدی، پایداری و دقت روش در پاسخ به پرسش‌های پژوهش مطرح‌شده در فصل اول است.


\section{محیط آزمایشگاهی، داده‌ها و معیارهای ارزیابی}\label{sec:5-1}

برای اطمینان از تکرارپذیری علمی نتایج، مشخصات بستر آزمایشگاهی، شیوه‌ی گردآوری یا تقسیم داده‌ها و معیارهای کمی سنجش باید با وضوح بیان شوند:
\begin{itemize}
    \item \textbf{بستر آزمون و مجموعه‌داده‌ها}: معرفی داده‌ها یا شرایط شبیه‌سازی، شامل تعداد نمونه‌ها، ابعاد ویژگی‌ها و شیوه‌ی تقسیم به بخش‌های آموزش، اعتبارسنجی و آزمون.
    \vspace{-0.25cm}
    \item \textbf{روش‌های خط‌پایه}: انتخاب روش‌های معتبر و متداول در ادبیات برای ایجاد مقایسه‌ای منصفانه.
    \vspace{-0.25cm}
    \item \textbf{معیارهای ارزیابی}: انتخاب سنجه‌های استاندارد متناسب با مسئله (نظیر دقت یا زمان پردازش).
\end{itemize}


\section{نتایج مقایسه‌ای و تحلیل تجربی}\label{sec:5-2}

در جدول‌های ارائه‌ی نتایج تجربی، مقادیر برتر معمولاً با قلم ضخیم (\textbf{\lr{Bold}}) یا با خط زیرین (\underline{\lr{Underline}}) مشخص می‌شوند. همچنین تکرار آزمایش‌ها در چندین دور مستقل و گزارش میانگین به همراه انحراف معیار، اعتبار آماری داده‌ها را نشان می‌دهد.

در جدول~\ref{tab:chap5_main_results}، نمونه‌ای استاندارد از جدول ارزیابی عملکرد روش پیشنهادی در مقایسه با روش‌های خط‌پایه بر روی سنجه‌های کمی ارائه شده است.

\begin{table}[!ht]
    \centering
    \small
    \setlength{\aboverulesep}{0pt}
    \setlength{\belowrulesep}{0pt}
    \setlength{\tabcolsep}{10pt}
    \caption{مقایسه‌ی روش پیشنهادی با خطوط‌پایه (میانگین $\pm$ انحراف معیار در ۵ تکرار)}\label{tab:chap5_main_results}
    \vspace{0.2cm}
    \begin{tabularx}{\textwidth}{l c c c}
        \toprule
        \textbf{روش ارزیابی‌شده} & \textbf{دقت عملکرد (\%)} & \textbf{خطای میانگین} & \textbf{زمان پردازش (ثانیه)} \\
        \midrule
        روش خط‌پایه‌ی اول (\lr{Baseline A}) & ۷۶٫۲۰ $\pm$ ۰٫۴۵ & ۰٫۲۳۸ $\pm$ ۰٫۰۰۸ & \textbf{۱٫۲} \\
        روش خط‌پایه‌ی دوم (\lr{Baseline B}) & ۸۳٫۱۵ $\pm$ ۰٫۳۰ & ۰٫۱۶۸ $\pm$ ۰٫۰۰۵ & ۲٫۵ \\
        روش خط‌پایه‌ی سوم (\lr{Baseline C}) & \underline{۸۸٫۷۰ $\pm$ ۰٫۲۵} & \underline{۰٫۱۱۳ $\pm$ ۰٫۰۰۴} & ۱۲٫۴ \\
        \midrule
        \textbf{روش پیشنهادی این پایان‌نامه} & \textbf{۹۲٫۴۰ $\pm$ ۰٫۱۸} & \textbf{۰٫۰۷۶ $\pm$ ۰٫۰۰۲} & \underline{۳٫۸} \\
        \bottomrule
    \end{tabularx}
\end{table}
\vspace{-0.5cm}


\section{تحلیل حساسیت و مطالعه‌ی اثر مؤلفه‌ها}\label{sec:5-3}

صرف ارائه‌ی جدول مقایسه‌ای برای یک پایان‌نامه‌ی پژوهشی کافی نیست؛ بلکه باید رفتار روش در شرایط گوناگون مورد موشکافی قرار گیرد. ارائه‌ی مطالعه‌ی اثر مؤلفه‌ها مشخص می‌سازد که هر یک از اجزای نوآوری پیشنهادی چه نقشی در بهبود نتایج داشته‌اند و عدم حضور هر بخش چه افتی در دقت یا کارایی ایجاد می‌کند. همچنین تحلیل حساسیت نسبت به تغییر ابرپارامترها، محدوده‌ی عملکرد ایمن سامانه را آشکار می‌سازد.


\section{بحث و بررسی تحلیلی یافته‌ها}\label{sec:5-4}

در پایان فصل آزمایش‌ها، نویسنده باید دلایل برتری یا ضعف‌های احتمالی روش را تبیین کند. صداقت علمی ایجاب می‌کند که شرایط شکست مدل یا سناریوهایی که روش در آن‌ها عملکرد ایده‌آلی ندارد نیز به صراحت بیان و ریشه‌یابی شوند. این تحلیل نقادانه نشان‌دهنده‌ی درک عمیق دانشجو از رفتار سیستم و زمینه‌ساز تدوین بخش محدودیت‌ها در فصل بعد است.
